\chapter{Description des Modules}

\section{ms-auth -- Service d'authentification}

\subsection{Responsabilités}
\begin{itemize}[itemsep=3pt]
  \item Gestion du cycle de vie des comptes utilisateurs (inscription, connexion, désactivation).
  \item Émission et validation des tokens JWT (accès 15 min + refresh 7 jours).
  \item Protection contre les attaques par force brute (\texttt{LoginRateLimiter}).
  \item Vérification email par lien tokenisé.
  \item Traçabilité de toutes les actions sensibles (audit logs).
\end{itemize}

\subsection{Endpoints principaux}

\begin{table}[H]
\centering
\renewcommand{\arraystretch}{1.3}
\begin{tabular}{|l|l|l|l|}
\hline
\rowcolor{tabhead}\color{white}\textbf{Méthode} &
\color{white}\textbf{URL} &
\color{white}\textbf{Accès} &
\color{white}\textbf{Description} \\
\hline
POST & \code{/api/v1/auth/register} & Public & Inscription patient \\
\rowcolor{roweven}
POST & \code{/api/v1/auth/login} & Public & Connexion + JWT \\
POST & \code{/api/v1/auth/refresh} & Public & Rafraîchir le token \\
\rowcolor{roweven}
POST & \code{/api/v1/auth/forgot-password} & Public & Demande reset mdp \\
POST & \code{/api/v1/auth/reset-password} & Public & Réinitialiser mdp \\
\rowcolor{roweven}
POST & \code{/api/v1/auth/change-password} & Connecté & Changer mdp \\
GET  & \code{/api/v1/auth/me} & Connecté & Profil courant \\
\rowcolor{roweven}
POST & \code{/api/v1/admin/personnel} & ADMIN & Créer compte personnel \\
\hline
\end{tabular}
\caption{Endpoints ms-auth}
\end{table}

\section{ms-patient-personnel -- Dossier médical}

\subsection{Responsabilités}
\begin{itemize}[itemsep=3pt]
  \item Gestion complète des dossiers médicaux électroniques.
  \item Portail patient : consultation, messagerie, documents, QR Code, export PDF.
  \item Dashboard direction : statistiques, accès tous patients, export PDF individuel.
  \item Proxy vers MedicalAppointments pour l'affichage des RDV.
\end{itemize}

\subsection{Structure du dossier médical}

\begin{table}[H]
\centering
\renewcommand{\arraystretch}{1.3}
\begin{tabular}{|l|l|}
\hline
\rowcolor{tabhead}\color{white}\textbf{Section} &
\color{white}\textbf{Contenu} \\
\hline
Antécédents & Médicaux, chirurgicaux, familiaux, allergies (avec sévérité) \\
\rowcolor{roweven}
Consultations & Date, médecin, motif, diagnostic, traitement, constantes vitales \\
Ordonnances & Médicaments, dosages, durées, instructions \\
\rowcolor{roweven}
Analyses labo & Type, laboratoire, statut, résultats \\
Radiologies & Type examen, description, conclusion \\
\rowcolor{roweven}
Hospitalisations & Service, dates entrée/sortie, motif \\
Certificats médicaux & Type, contenu, date émission \\
\rowcolor{roweven}
Documents uploadés & Ordonnances scannées, radios, analyses (PDF/images) \\
\hline
\end{tabular}
\caption{Structure du dossier médical électronique}
\end{table}

\section{MedicalAppointments -- Prise de rendez-vous}

\subsection{Responsabilités}
\begin{itemize}[itemsep=3pt]
  \item Gestion des créneaux horaires par les médecins (création unitaire ou en masse, blocage).
  \item Réservation de rendez-vous pour les patients authentifiés \emph{ou anonymes}.
  \item Annulation avec token anonyme (pour les patients sans compte).
  \item Liste d'attente avec notification email automatique à la libération d'un créneau.
  \item Notifications temps réel via SignalR (\code{/appointmenthub}).
\end{itemize}

\subsection{Statuts d'un rendez-vous}

\begin{table}[H]
\centering
\renewcommand{\arraystretch}{1.3}
\begin{tabular}{|l|l|l|}
\hline
\rowcolor{tabhead}\color{white}\textbf{Statut} &
\color{white}\textbf{Description} &
\color{white}\textbf{Transition depuis} \\
\hline
Confirmed & RDV confirmé par défaut à la réservation & (initial) \\
\rowcolor{roweven}
Pending & En attente de confirmation manuelle & Confirmed \\
CancelledByPatient & Annulé par le patient & Confirmed, Pending \\
\rowcolor{roweven}
CancelledByDoctor & Annulé par le médecin & Confirmed, Pending \\
CancelledBySecretary & Annulé par la secrétaire & Confirmed, Pending \\
\rowcolor{roweven}
Completed & Consultation effectuée & Confirmed \\
NoShow & Patient absent & Confirmed \\
\hline
\end{tabular}
\caption{Cycle de vie d'un rendez-vous}
\end{table}

\subsection{Règle de pénalité patient}
Si un patient annule ou est absent plus de 3 fois, le système enregistre
une date de \texttt{PenaltyUntil} qui bloque temporairement ses réservations.
Le compteur \texttt{CancelCount} est incrémenté à chaque annulation.

\section{ms-personnel -- Gestion RH}

\subsection{Responsabilités}
\begin{itemize}[itemsep=3pt]
  \item Gestion des profils du personnel médical et paramédical.
  \item Planifications des gardes et disponibilités.
  \item Gestion des absences (congés, maladie, formation).
  \item Organisation en équipes avec responsable.
  \item Demandes de modification de planning.
\end{itemize}

\subsection{Types de personnel gérés}

\begin{itemize}
  \item Médecins (avec spécialité, créneaux de consultation)
  \item Infirmiers
  \item Aides-soignants
  \item Brancardiers
  \item Secrétaires médicales
  \item Directeurs de service
\end{itemize}

\section{medsys-web -- Interface utilisateur}

\subsection{Tableaux de bord par rôle}

\begin{table}[H]
\centering
\renewcommand{\arraystretch}{1.3}
\begin{tabular}{|l|l|}
\hline
\rowcolor{tabhead}\color{white}\textbf{Rôle} &
\color{white}\textbf{Fonctionnalités accessibles} \\
\hline
PATIENT & Accueil, Dossier médical, Prendre RDV, Mes RDV,
  Documents, Messagerie, QR/PDF, Profil \\
\rowcolor{roweven}
MÉDECIN / PERSONNEL & Liste patients, Dossier complet par patient \\
DIRECTEUR & Vue d'ensemble (KPIs), Patients, Médecins,
  Comptes, RDV \& Planning, Stats RDV \\
\rowcolor{roweven}
ADMIN & Gestion des comptes (créer, activer/désactiver) \\
\hline
\end{tabular}
\caption{Fonctionnalités par rôle dans medsys-web}
\end{table}

\subsection{Flux d'authentification frontend}

\begin{enumerate}[itemsep=3pt]
  \item Le patient se connecte via \textbf{LandingPage} (split login).
  \item Le token JWT est stocké en \code{sessionStorage} sous la clé \code{medsys\_token}.
  \item \textbf{AuthContext} décode le token et expose le rôle à tous les composants.
  \item React Router redirige vers le dashboard approprié selon le rôle.
  \item À l'expiration du token (15 min), un refresh automatique est tenté.
\end{enumerate}

\subsection{Intégration MedicalAppointments}
La section \textbf{Prendre RDV} du dashboard patient utilise directement
l'API MedicalAppointments via le proxy Vite. Le flux est :

\begin{enumerate}[itemsep=3pt]
  \item \textbf{Étape 1} -- Sélection du service médical (\code{GET /api/services}).
  \item \textbf{Étape 2} -- Sélection du médecin (\code{GET /api/services/\{id\}/doctors}).
  \item \textbf{Étape 3} -- Vue semaine avec créneaux (\code{GET /api/slots?doctorId=\&weekStart=}).
  \item \textbf{Confirmation} -- Modal de confirmation avec motif, réservation
    (\code{POST /api/appointments}), affichage du token d'annulation si anonyme.
  \item \textbf{Liste d'attente} -- Si semaine complète : \code{POST /api/waiting-list}.
\end{enumerate}
